\documentclass{article}
\usepackage{tcee-repaso}
\tcetema[El Fondo Monetario Internacional]{3.B.22}{El Fondo Monetario Internacional. Estructura y políticas. La prevención y solución de crisis.}
\tcefecha{13/11/2026}
\begin{document}
\cabecerarepaso
\begin{columnas}

% ==========================================================================
\bloque{Esquema del cante}
\esquemacante{%
\apartadocante{Introducción: «más vale prevenir que curar»}{2,5}
\apartadocante{I. Estructura}{8,5}
\subcante{1. Evolución histórica}{2,5}
\subcante{2. Marco jurídico}{1,5}
\subcante{3. Marco institucional}{4,5}
\apartadocante{II. Políticas}{8,5}
\subcante{1. Supervisión}{3}
\subcante{2. Asistencia técnica}{1}
\subcante{3. Política financiera}{4,5}
\apartadocante{III. Prevención y solución de crisis}{9,5}
\subcante{1. Idea: el FMI como prestamista internacional}{2}
\subcante{2. La prevención}{2}
\subcante{3. La solución: los documentos del programa}{3,5}
\subcante{4. Consecuencias de los programas}{1}
\subcante{5. Valoración}{1}
\apartadocante{Conclusión}{1}}
\emph{Gráficos del cante}: los votos (\hyperref[rep:I]{I}), el crédito vivo (\hyperref[rep:II]{II}) y las 3 regiones de la financiación catalítica (\hyperref[rep:III]{III}); fórmulas: votos, techo de crédito y métrica ARA.

% ==========================================================================
\bloque[2,5]{Introducción}
\begin{itemize}
\item \textbf{Enganche}: «Más vale prevenir que curar». Argentina, abril de 2025: 20.000 M\,\$ (479\,\% de la cuota); debe $\approx$42.000 M de DEG, casi la mitad de la cuenta general. ¿Puede el FMI evitar las crisis o solo curarlas, y a qué precio?
\item \textbf{Relevancia}: 191 miembros; capacidad $\approx$1 billón de \$, centro de la red de seguridad; España miembro desde 1958, Plan de Estabilización de 1959 \vertema{4.A.27}, banca 2012 \vertema{4.A.28}; informes de referencia (WEO, GFSR, art.~IV).
\item \textbf{Contextualización}: \concepto{organización internacional} (Convenio Constitutivo, Bretton Woods, julio de 1944, 44 países), organismo especializado de la ONU. \emph{Fondo} (recursos por cuotas), \emph{monetario} (de los bancos centrales; refuerza reservas, no financia proyectos) e \emph{internacional}. Art.~I: cooperación, comercio, estabilidad cambiaria, pagos multilaterales, prestar para corregir la BP sin medidas perniciosas. \concepto{Crisis de BP} \vertema{3.B.17}; \concepto{red de seguridad financiera mundial}: reservas, \emph{swaps}, fondos regionales (MEDE, Chiang Mai) y FMI.
\item \textbf{Problemática}: ¿cómo se gobierna y se financia? ¿Qué hace? ¿Previene o solo cura? ¿A qué coste (riesgo moral, ajuste, legitimidad)?
\item \textbf{Estructura}: \hyperref[rep:I]{I} qué es; \hyperref[rep:II]{II} qué hace; \hyperref[rep:III]{III} para qué: prevenir y resolver.
\end{itemize}

% ==========================================================================
\bloque[8,5]{I. Estructura\label{rep:I}}
\begin{ideaclave}
Nació para vigilar tipos fijos y se reinventó como supervisor y prestamista de crisis. Gobierno por accionistas: quien más aporta, más vota; Estados Unidos tiene el veto (85\,\%).
\end{ideaclave}

\sub{I.1. Evolución histórica\label{rep:I.1}}
\begin{itemize}
\item \emph{Bretton Woods (1944-1973)}: tipos fijos pero ajustables (dólar-oro a 35\,\$); planes \autor{Keynes} (\emph{bancor}, ajuste simétrico) vs. \autor{White} (fondo por cuotas): gana \autor{White}. Paridades: cambio solo con \concepto{desequilibrio fundamental} ($>$10\,\% con conformidad). 1952 Stand-By, condicionalidad; 1962 AGP; 1969 DEG (\autor{Triffin}); 1971-1973 fin \vertema{3.B.20}.
\item \emph{1973-2007}: 2.ª enmienda (1976-1978): libre régimen salvo oro, \emph{supervisión firme}. Deuda de 1982, préstamos de favor (1986-1987), crisis emergentes 1994-2002 \vertema{3.B.21}; 2002 \concepto{acceso excepcional}; 2003-2007 devoluciones, mínimo de crédito.
\item \emph{2008-2019}: Londres 2009 (recursos ×3, 750.000 M\,\$); LCF; reforma de 2010 (cuotas ×2, $>$6\,\% a emergentes, en vigor 2016); troika; Grecia 2010 ($>$3.200\,\% de la cuota) con \concepto{excepción sistémica}, suprimida en 2016.
\item \emph{2020-2026}: urgencias, FFACC, DEG 2021 (456.000 M); SRS (2022); Mesa Redonda de la deuda (2023); Ucrania; Bessent (2025): «que el FMI vuelva a ser el FMI».
\end{itemize}

\graficorepaso[\linewidth]{cronologia-fmi}{4 etapas: guardián de las paridades, supervisor y bombero, regreso con la crisis global y el euro, pandemia y después}

\sub{I.2. Marco jurídico\label{rep:I.2}}
\begin{itemize}
\item Convenio en vigor 27-12-1945 (opera desde 1-3-1947); 7 enmiendas (3/5 de miembros con el 85\,\%). Peso de los fines: paridades (BW) $\to$ financiación y cooperación.
\item Obligaciones: \emph{art.~IV} (no manipular, supervisión); \emph{art.~VIII} (sin restricciones a pagos corrientes; transición art.~XIV); \emph{art.~VI} (controles de capital permitidos; Visión institucional 2012 y 2022 \vertema{3.B.15}); datos (sanciones: Argentina 2013).
\end{itemize}

\sub{I.3. Marco institucional\label{rep:I.3}}
\begin{itemize}
\item \textbf{Junta de Gobernadores} (191; reservado: cuotas, DEG, miembros, enmiendas, 85\,\%) $\to$ \textbf{Directorio} (25 sillas desde nov. 2024; consenso; España con México y Centroamérica) $\to$ \textbf{Director gerente} (\autor{Georgieva}, 2.º mandato; europea por tradición) y $\approx$3.100 empleados. Asesoran el \emph{CMFI} (\autor{Aljadaan}) y el Comité para el Desarrollo; evalúa la OEI (2011: no vio 2008; 2016: troika).
\end{itemize}
\graficorepaso[\linewidth]{gobierno-fmi}{Gobierno: Junta $\to$ Directorio $\to$ Gerencia; asesoran CMFI y Comité para el Desarrollo; evalúa la OEI}
\begin{itemize}
\item \textbf{Cuotas}: suscripción (25\,\% en DEG o divisas), votos, acceso, DEG. Total 476.000 M de DEG; 16.ª revisión (+50\,\%, 715.700 M) sin vigor: 76,7\,\% de consentimientos, falta EE.~UU.; 17.ª con los Principios de Diriyah.
\begin{formula}
\centering $\text{Votos}_i=\text{básicos (5,502~\%)}+\dfrac{\text{Cuota}_i}{100.000\ \text{DEG}}$\\[2pt]
$\text{CQS}_i=(0{,}5\,\text{PIB}+0{,}3\,\text{Ap}+0{,}15\,\text{Var}+0{,}05\,\text{Res})^{0{,}95}$
\end{formula}
\end{itemize}
\graficorepaso[\linewidth]{cuotas-votos}{Votos: EE.~UU. 16,49\,\% $>$ 15\,\% (veto); Japón 6,14; China 6,08; España 1,92 (13.º)}
\begin{itemize}
\item PIB 60\,\% mercado/40\,\% PPA; apertura; variabilidad; reservas. Avanzadas: 61,6\,\% de las cuotas.
\item \textbf{Recursos prestados}: NAP (364.000 M de DEG, 2026-2030) y 39 bilaterales; nunca ha emitido bonos.
\item \textbf{DEG}: activo de reserva y unidad de cuenta; cesta 2022: \$ 43,38, € 29,31, CNY 12,28, ¥ 7,59, £ 7,44 (hasta jul. 2027); tipo 3,119\,\% (nov. 2026); 660.800 M asignados.
\end{itemize}
\graficorepaso[\linewidth]{cesta-deg}{Cesta del DEG: sube el dólar y el renminbi, baja el euro (2015 $\to$ 2022)}
\begin{itemize}
\item \textbf{Balance}: CRG (cuotas, préstamos, reservas / monedas, crédito, oro 2.814 t) y fondos fiduciarios FFCLP y FFRS. \concepto{Acreedor preferente}; saldos precautorios 26.300 M.
\end{itemize}

% ==========================================================================
\bloque[8,5]{II. Políticas\label{rep:II}}
\begin{ideaclave}
Vigila a todos (solo persuade), ayuda a quien lo pide y presta con condiciones a quien lo necesita.
\end{ideaclave}

\sub{II.1. Supervisión\label{rep:II.1}}
\begin{itemize}
\item Base: art.~IV y \concepto{Decisión de Supervisión Integrada} (2012, contagios); revisiones de 2021 y 2026.
\item \emph{Bilateral}: consultas del art.~IV (anuales); \concepto{FSAP} (1999, con el BM; obligatorio desde 2010 cada 5 años: 47 jurisdicciones; España 2006, 2012, 2017, 2024); supervisión cambiaria \emph{de facto} vs. \emph{de iure}, \concepto{miedo a flotar} (\autor{Calvo}-\autor{Reinhart}); no sanciona la manipulación (China).
\item \emph{Multilateral}: WEO, GFSR y Monitor fiscal (semestrales); ESR (anual); alerta temprana con el FSB; G-20. \emph{Regional}: perspectivas regionales.
\item \emph{Estadística}: NEDD y SGDD-m; MBP7 (2025; aplicación 2029-2030) \vertema{3.B.11}.
\end{itemize}

\sub{II.2. Asistencia técnica\label{rep:II.2}}
\begin{itemize}
\item A petición: monetaria y cambiaria, financiera, tributaria, gasto, deuda, estadística; asesoramiento y formación. España 2012 \vertema{4.A.18}.
\end{itemize}

\sub{II.3. Política financiera\label{rep:II.3}}
\begin{itemize}
\item \emph{Justificación}: crisis recurrentes (\autor{Reinhart}-\autor{Rogoff}); información (\autor{Diamond}-\autor{Dybvig}, Nobel 2022); bien público (\autor{Stiglitz}, Nobel 2001). Pero \concepto{riesgo moral}.
\item \emph{Principios}: universalidad, temporalidad y garantías, sin lucro con incentivos, \concepto{condicionalidad} (tramos; techo de CI, suelo de RIN, déficit; referencias estructurales). Técnicamente, \emph{compras} y \emph{recompras}.
\item \emph{CRG}: Stand-By (corto plazo; 3¼-5 años), SAF (estructural; 4½-10), LCF y LPL (prevención), LLCP (se extingue en abr. 2027), IFR (urgencia). \emph{FFCLP}: SCA, SCS, SCR. \emph{FFRS}: SRS (20 años). Sin dinero: ICP, IAPE.
\end{itemize}
{\footnotesize\centering
\begin{tabular}{@{}>{\raggedright\arraybackslash}p{0.2\linewidth}>{\raggedright\arraybackslash}p{0.4\linewidth}>{\raggedright\arraybackslash}p{0.3\linewidth}@{}}
\toprule
\textbf{Instrumento} & \textbf{Para qué} & \textbf{Duración / reembolso}\\
\midrule
Stand-By & BP a corto plazo & $\leq$3 años / 3¼-5\\
SAF & estructural, reformas & $\leq$4 años / 4½-10\\
LCF & prevención, sin condiciones \emph{ex post} & 1-2 años / 3¼-5\\
LPL & prevención con vulnerabilidades & 6 m-2 años / 3¼-5\\
IFR, SCR & urgencias & único / 3¼-5, 5½-10\\
SCA, SCS & renta baja (FFCLP) & 3-5, 1-3 años\\
SRS & clima y pandemias (FFRS) & $\geq$18 m / 20 (10½ gracia)\\
\bottomrule
\end{tabular}\par}
\begin{itemize}
\item \emph{Acceso}: 200\,\% anual, 600\,\% acumulado; por encima, \concepto{acceso excepcional} (4 criterios: necesidad, deuda sostenible con alta probabilidad, vuelta a mercados, capacidad de ajuste).
\item \emph{Coste}: tipo del DEG + 60 pb; recargos de 200 pb ($>$300\,\% de la cuota) y +75 pb ($>$3 años); reforma de 2024: $-$36\,\%.
\end{itemize}
\graficorepaso[\linewidth]{credito-fmi}{Crédito vivo (miles de M de DEG): oleadas de 1982, 1995-2003, euro y pandemia; mínimo en 2007; 122.700 M en oct. 2026, Argentina 46\,\% de la CRG}

% ==========================================================================
\bloque[9,5]{III. Prevención y solución de crisis\label{rep:III}}
\begin{ideaclave}
El FMI evita crisis de liquidez si el país es solvente y llega a tiempo; previene con supervisión, reservas y líneas precautorias, y resuelve con financiación, ajuste y, si la deuda no es sostenible, reestructuración. Precio: riesgo moral y ajustes dolorosos.
\end{ideaclave}

\sub{III.1. Idea: financiación catalítica\label{rep:III.1}}
\begin{itemize}
\item Prestamista de última instancia (\autor{Bagehot}); \autor{Fischer} (1999): el FMI, más gestor de crisis (no crea dinero, sin garantías).
\item \emph{Modelo} (\autor{Corsetti}-\autor{Guimarães}-\autor{Roubini}, 2006; \autor{Morris}-\autor{Shin}): deuda $D$, reservas $R$, inversión $\theta$ (liquidación a $\psi<1$), préstamo $L$. Insolvente si $\theta<\theta_I=D-R-L$; a salvo si $\theta\geq\theta_S=(D-R-L)/\psi$; entre ambos, pánico posible (juegos globales: umbral único $\theta^*$).
\begin{formula}
\centering $\theta_S-\theta_I=(D-R-L)\left(\dfrac{1}{\psi}-1\right)$
\end{formula}
\item Implicaciones: evita crisis de liquidez; no salva a un insolvente (reestructurar); tamaño y rapidez; el riesgo moral no es inevitable. Riesgo moral del acreedor (\autor{Jeanne}-\autor{Zettelmeyer}, «teorema de \autor{Mussa}»); evidencia: \autor{Mody}-\autor{Saravia} (2006), rescate de Rusia 1998.
\end{itemize}
\graficorepaso[\linewidth]{catalitica}{Con $L$, los umbrales se mueven a la izquierda y la zona de pánico se estrecha}

\sub{III.2. La prevención\label{rep:III.2}}
\begin{itemize}
\item Supervisión; autoseguro con reservas; seguro común (líneas precautorias).
\item \textbf{Adecuación de reservas} (tema de referencia): 3 meses de importaciones; 100\,\% de la deuda a corto (\autor{Greenspan}-\autor{Guidotti}); $\approx$20\,\% de M2; métrica ARA:
\begin{formula}
\centering $\text{ARA}=a_X X+a_M\mathrm{M2}+a_D\mathrm{DCP}+a_P\mathrm{OP}$\\[1pt]
\footnotesize fijo 10/10/30/20~\%; flexible 5/5/30/15~\%; adecuado 100-150~\%
\end{formula}
\item Coste $\approx$1\,\% del PIB (\autor{Rodrik}, 2006); óptimo (\autor{Jeanne}-\autor{Rancière}, 2011); en 2008 protegieron (\autor{Obstfeld}-\autor{Shambaugh}-\autor{Taylor}; \autor{Frankel}-\autor{Saravelos}). Límite: pesos del pasado, sin coste ni \emph{swaps}.
\item \textbf{LCF} (2009, sin condiciones \emph{ex post}: México, Colombia, Polonia, Chile, Perú, Marruecos), \textbf{LPL}, LLCP. \emph{Estigma}; la red crece fuera del FMI, desigual.
\end{itemize}

\sub{III.3. La solución: documentos del programa\label{rep:III.3}}
\begin{itemize}
\item \textbf{Programación financiera} (\autor{Polak}, 1957): tipo fijo; $\Delta M=\Delta R+\Delta\mathrm{CI}$; $M=k\,PY$; $\Delta R=X-m\,PY+\Delta F$. Un $\Delta$CI sube la renta $1/(k+m)$ y baja $R$ en $m/(k+m)$; a largo plazo $\Delta R=-\Delta\mathrm{CI}$ \vertema{3.B.12}.
\begin{formula}
\centering $\Delta\mathrm{CI}^{\max}=k\,\Delta(PY)-\Delta R^*$
\end{formula}
\item[] Techo de CI y suelo de RIN; ajuste de demanda; 4 cuentas coherentes. Extensiones: modelo fusionado FMI-BM; con metas de inflación, cláusula de consulta; presentación analítica y brecha de financiación. Valoración: coherencia contable, sin expectativas ni enfoque intertemporal.
\item \textbf{Art.~IV}: real, exterior ($\mathrm{CC}=\mathrm{RNDB}-A=S-I=(S_P-I_P)+(T-G)$; reducir, desviar el gasto, ampliar la oferta), monetario, público.
\item \textbf{DSA}: MAC SRDSF (2021: riesgo bajo/moderado/alto) y LIC-DSF; filtro: alta probabilidad $\to$ presta; no alta $\to$ alargar plazos; insostenible $\to$ reestructurar \vertema{3.B.29}.
\item \textbf{Sector financiero}: activos, solvencia, pruebas de resistencia; liquidez, recapitalización, bancos puente y malos, \emph{bail-in}, controles («corralito», Ecuador 1999).
\item \textbf{Reestructuración}: SDRM (\autor{Krueger}, 2001-2003, abandonado); CAC (México 2003; reforzadas 2014); préstamo con atrasos; Marco Común (Chad, Zambia, Ghana, Etiopía) y Mesa Redonda; 3 pilares (2024).
\end{itemize}

\sub{III.4. Consecuencias\label{rep:III.4}}
\begin{itemize}
\item Corto plazo: $\downarrow$crecimiento, $\uparrow$deuda/PIB, $\downarrow$inflación, $\uparrow$CC, depreciación que se estabiliza. Medio (6 meses-2 años): vuelve el crecimiento, $\uparrow$reservas, $\downarrow$déficit.
\item Datos: $\downarrow$crecimiento (\autor{Przeworski}-\autor{Vreeland}; \autor{Barro}-\autor{Lee}), menos si se cumple (\autor{Dreher}); multiplicadores infravalorados (\autor{Blanchard}-\autor{Leigh}, 2013). Asia 1997 (\autor{Stiglitz}, \autor{Krugman}); Grecia (evaluación \emph{ex post} 2013); Argentina (2018, 2025).
\end{itemize}

\sub{III.5. Valoración\label{rep:III.5}}
\begin{itemize}
\item \emph{A favor}: financia con el mercado cerrado, credibilidad (catálisis), frena contagios, bien público.
\item \emph{En contra}: riesgo moral del acreedor (tesobonos, 1995); coste social y errores (apertura antes de 1997, austeridad); supervisión sin dientes; legitimidad (cuotas, veto, reparto UE-EE.~UU.); recursos limitados.
\end{itemize}

% ==========================================================================
\bloque[1]{Conclusión}
\begin{itemize}
\item \textbf{Recapitulación}: «En conclusión,» fondo monetario de accionistas reinventado como supervisor y prestamista; 3 políticas para prevenir y resolver; cura la iliquidez, no la insolvencia.
\item \textbf{Relevancia}: fragmentación, deuda de renta baja y media, clima; centro de la red, cada vez más compartido.
\item \textbf{Extensiones}: \vertema{3.B.20} \vertema{3.B.21} \vertema{3.B.17} \vertema{3.B.29} \vertema{3.B.30} \vertema{3.B.19} \vertema{4.A.21}.
\item \textbf{Opinión}: ha aprendido (controles, gasto social, deuda); falta la legitimidad de las cuotas.
\item \textbf{Idea final}: más vale prevenir que curar, pero alguien tiene que curar: a tiempo, al solvente y sin quitarle las ganas de cuidarse.
\end{itemize}

% ==========================================================================
\bloque{Autores y fechas}
\hito{1873}{\autor{Bagehot}: prestamista de última instancia}
\hito{1944}{Bretton Woods: \autor{Keynes} vs. \autor{White}}
\hito{1945}{Convenio en vigor (27-12); opera desde 1947}
\hito{1952}{Acuerdos Stand-By}
\hito{1958}{Entra España}
\hito{1962}{AGP del G-10}
\hito{1957}{\autor{Polak}: programación financiera}
\hito{1969}{1.ª enmienda: DEG}
\hito{1976}{Jamaica: 2.ª enmienda}
\hito{1982}{Crisis de la deuda; \autor{Volcker}}
\hito{1983}{\autor{Diamond}-\autor{Dybvig}: pánicos}
\hito{1986}{Servicio de Ajuste Estructural (1987, reforzado)}
\hito{1999}{\autor{Fischer}: prestamista internacional; \autor{Greenspan}-\autor{Guidotti}; FSAP}
\hito{1998}{NAP}
\hito{2001}{OEI; \autor{Krueger}: SDRM}
\hito{2002}{\autor{Calvo}-\autor{Reinhart}: miedo a flotar; acceso excepcional; \autor{Stiglitz}}
\hito{2005}{\autor{Jeanne}-\autor{Zettelmeyer}: teorema de \autor{Mussa}}
\hito{2006}{CGR y \autor{Morris}-\autor{Shin}: catálisis; \autor{Rodrik}: coste de las reservas}
\hito{2009}{Londres: recursos ×3; LCF}
\hito{2011}{Métrica ARA; \autor{Jeanne}-\autor{Rancière}}
\hito{2012}{Supervisión integrada; Visión institucional}
\hito{2013}{\autor{Blanchard}-\autor{Leigh}: multiplicadores}
\hito{2016}{Fin de la excepción sistémica; reforma de 2010 en vigor}
\hito{2020}{Préstamos de urgencia; FFACC; LLCP}
\hito{2021}{DEG 456.000 M; MAC SRDSF}
\hito{2022}{SRS (clima); revisión de la Visión institucional}
\hito{2023}{16.ª revisión; Mesa Redonda}
\hito{2024}{25.ª silla; reforma de recargos; 3 pilares}
\hito{2025}{SAF Argentina; MBP7; Bessent}
\hito{2026}{SAF Ucrania; revisiones de supervisión, condicionalidad y FSAP; Bangkok}

% ==========================================================================
\bloque{Datos clave}
\begin{itemize}
\item Cargos: \autor{Georgieva} (gerente), \autor{Katz} (primer subdirector), \autor{Tenreyro} (economista jefe), \autor{Aljadaan} (CMFI).
\item Miembros \dato{191}{FMI}; Directorio \dato{25 sillas}{nov. 2024}; personal \dato{3.118}{abr. 2026}.
\item Cuotas \dato{476.000 M de DEG}{FMI}; 16.ª \dato{715.700 M, 76,7\,\% de consentimientos}{CRS, 2026}.
\item Votos: EE.~UU. \dato{16,49\,\%}{FMI}; España \dato{1,92\,\% (cuota 2\,\%)}{FMI}.
\item NAP \dato{364.000 M de DEG}{FMI}; capacidad \dato{$\approx$1 billón \$}{FMI}.
\item Crédito vivo \dato{122.700 M de DEG}{FMI, oct. 2026}; Argentina \dato{46\,\% de la CRG}{FMI}.
\item Oro \dato{2.814 t}{FMI}; saldos precautorios \dato{26.300 M de DEG}{abr. 2026}.
\item Asignaciones de DEG \dato{660.800 M; 456.000 en 2021}{FMI}.
\item Argentina: SAF \dato{20.000 M\,\$, 479\,\% de la cuota}{abr. 2025}; Ucrania: SAF \dato{8.100 M\,\$}{feb. 2026}.
\item DEG: \dato{1 DEG = 1,347\,\$; 3,119\,\%}{FMI, nov. 2026}.
\end{itemize}

% ==========================================================================
\bloque{Preguntas del tribunal}
\begin{enumerate}
\item \textbf{¿Es el FMI un prestamista de última instancia?} No del todo: no crea dinero ni toma garantías; gestor de crisis (\autor{Fischer}).
\item \textbf{¿Por qué tiene veto EE.~UU.?} Grandes decisiones con el 85\,\% y EE.~UU. tiene el 16,49\,\%.
\item \textbf{¿Qué es un DEG?} Activo de reserva y unidad de cuenta; derecho a obtener monedas; no es crédito contra el FMI.
\item \textbf{¿Por qué no entra en vigor la 16.ª revisión?} Falta el 85\,\%: el Congreso de EE.~UU.
\item \textbf{¿Cuándo presta por encima de los límites?} Acceso excepcional: 4 criterios; sin excepción sistémica desde 2016.
\item \textbf{¿Cuántas reservas bastan?} Greenspan-Guidotti y ARA (100-150\,\%); cuestan $\approx$1\,\% del PIB.
\item \textbf{¿Crea riesgo moral?} Del acreedor, sí; del deudor, poco (teorema de \autor{Mussa}); la catálisis puede aumentar el esfuerzo.
\item \textbf{¿Por qué fracasó Grecia?} Ajuste sin reestructurar, previsiones optimistas, multiplicadores infravalorados.
\item \textbf{¿Qué diferencia hay entre el FMI y el Banco Mundial?} El FMI refuerza reservas y corrige la BP a corto y medio plazo; el BM financia proyectos y desarrollo \vertema{3.B.30}.
\item \textbf{¿Qué es la condicionalidad y por qué existe?} Desembolsos por tramos contra criterios cuantitativos y referencias estructurales: protege los recursos y reduce el riesgo moral; revisión de 2026: ajuste adelantado si es crítico.
\item \textbf{¿Por qué los emergentes acumulan reservas?} Autoseguro tras 1997-1998, estigma de pedir al FMI y motivo mercantilista; problema colectivo.
\item \textbf{¿Qué es la programación financiera?} Marco de coherencia de 4 cuentas (\autor{Polak}): techo de CI = $k\Delta(PY)-\Delta R^*$.
\item \textbf{¿Qué hace el FMI en una reestructuración?} DSA (cuánto alivio), préstamo con atrasos, Marco Común y Mesa Redonda; CAC frente al SDRM.
\item \textbf{¿Qué es la red de seguridad financiera mundial?} Reservas, \emph{swaps}, fondos regionales y FMI; cobertura desigual.
\item \textbf{¿Qué supervisa el art.~IV y cómo?} Misión anual, informe al Directorio; 4 sectores (real, exterior, monetario, público); efectos de contagio desde 2012.
\item \textbf{¿Por qué se suprimió la excepción sistémica?} Permitía prestar a Grecia sin deuda sostenible con alta probabilidad; desde 2016, alargar vencimientos con ayuda de otros acreedores.
\item \textbf{¿Qué críticas hizo la OEI?} 2011: sesgo de grupo, no vio 2008; 2016: excepción sistémica y papel en la troika.
\item \textbf{¿Cómo pesa España en el FMI?} Cuota 9.535,5 M de DEG (2\,\%), 1,92\,\% de los votos, silla con México y Centroamérica; FSAP cada 5 años (2024).
\end{enumerate}

\begin{trampa}
\begin{itemize}
\item La 14.ª revisión (2010) entró en vigor en \emph{2016}; la 15.ª (2020) sin aumento; la 16.ª (2023) aún no está en vigor.
\item El FMI \emph{compra} y \emph{recompra}: el crédito es lo que tiene de la moneda del país por encima de su cuota.
\item Recargo: 200 pb por encima del \emph{300}\,\% de la cuota (antes 187,5\,\%); margen 60 pb (antes 100).
\item 85\,\% de los \emph{votos}, no de los miembros; las enmiendas exigen además 3/5 de los miembros.
\item El art.~VI \emph{permite} controles de capital; el VIII prohíbe restringir pagos \emph{corrientes}.
\item El DEG no es moneda; cesta de 5 (renminbi desde 2016); se revisa cada 5 años.
\item El PRGT (FFCLP) no está en la CRG; las cuotas: 25\,\% en reservas, 75\,\% en moneda propia.
\item La LCF no tiene condiciones \emph{ex post}; la LLCP se extingue en 2027.
\item El FSAP es con el Banco Mundial (1999), obligatorio desde 2010.
\item Directora gerente: europea; presidente del BM: estadounidense.
\item \autor{Polak}: tipo \emph{fijo}; a largo plazo todo el crédito se va en reservas.
\end{itemize}
\end{trampa}

\end{columnas}
\end{document}
